\section{Limitations and conclusion}
\label{sec:limitations}

\paragraph{Conclusion.}
We prove, in a serial linear-encoder model, that a contextual pathway
trained faster than the spectral encoder suppresses the encoder's update
at matched fit, that the fitted model fails under contextual reversal
under explicit finite-width conditions, and that readout normalization
removes the width dependence without necessarily removing the failure.
At matched training loss in the synthetic spectral--spatial model, slowing
the spatial head raises the linear accessibility of the spectral cue
while the trained classifier remains strongly context-reliant; retraining
a fresh head on decorrelated context recovers about $0.2$ more shifted
accuracy after the slow head than after the fast one, which yields no
more than a random encoder. With measured spectra, high accessibility
coexists with severe reliance, so a probe deficit is not necessary for
failure; with low accessibility the head's rate matters more at the
smaller constructed amplitude. The
paper's claim is that difference, at the stated fits, rates and
constructions, not a general law.

\paragraph{Limitations.}
The theorem has a linear encoder, noiseless cues and a fixed spectral skip
coefficient; the experiments add noisy cues and jointly learned routing
but keep a linear encoder and one head family each. The cues that show
the mechanism are constructed; natural $3\times3$ neighbourhoods, being
the same class as their centre, did not produce it, and larger receptive
fields are untested. Matched mean training loss is an analogue of the
theorem's matched margin, not the same rule; where a random encoder
already reads the spectral cue the failure is not the theorem's.
Real-spectrum seed variation is conditional on one patient split. On a
production model, scalar curvature is not the diagnostic used here and an
earlier freezing advantage was protocol-confounded
(Appendix~\ref{app:scope}). We test no mitigation on natural context;
what the experiments support is conditional: where the spectral cue is
accessible at the encoder output, readout retraining on
context-decorrelated data recovered it; where it is not, only encoder
retraining without informative context did.
